\section{宫廷的摇摆与北方的新行动者}

1100年哲宗无嗣而死，徽宗即位。端王被选择的具体过程材料有限，后出的预言与人物评语不能替代可核验的宫廷程序。1102年蔡京再入中枢，1104年元祐党人碑公开标列异议者；1106年朝廷又一度调整崇宁政策、宽释党人，1107年蔡京复相。四个节点见于《宋史》徽宗本纪、蔡京传和会要相关条目，足以说明人事与政策的反复，却不支持把所有晚期政务或全部危机归咎于一人。党籍、学校、财政和任官被连在一起，令朝廷的短期调整更难摆脱既有依附网络。

1115年完颜阿骨打建金、反辽，北方格局因此重组。《金史》太祖本纪、辽天祚帝本纪和宋徽宗本纪可以交叉确证建国与战争的基本序列；宋廷初获何种情报、何时形成判断，仍须细读使节材料。女真集团的兴起既来自辽统治下的压力，也依赖完颜联盟的组织能力，不宜写成宋朝外交所能任意操纵的机会。

1118年宋廷经海路和边境渠道接触金方，1120年达成海上之盟，共攻辽并商议燕云。宋、金交聘表与《三朝北盟会编》都保存谈判和盟约线索，但条款传本、岁币转换和燕云范围有异文。此盟改变了宋辽之间已延续多年的停战框架：宋希望借金军收复燕云，金则需要南方牵制辽军。它不是一份自动保证收复的契据，而是把宋的军费、作战能力和战后边界安排置于更严厉的检验之下。

\section{燕京的交接不是收复的终点}

1121年宋军在两浙平定方腊后，1122年童贯等依盟约攻辽失利；同年金军攻取辽南京燕京，辽朝北撤。宋、辽、金三方记载对军纪、抵抗和失利缘由各有立场，城内破坏与居民处境也不能由任一方的胜利叙述概括。可以把握的是，宋未能独力取得燕京，金因而掌握交接条件。

1123年金将燕京及部分地区交给宋，宋承担岁币与边防成本。名义上的移交不等于实际控制：范围、户口、城防和驻军须逐项核对。新前沿既脆弱又昂贵，宋金互信迅速下降。1123--1124年平州守将张觉降宋，宋廷接纳后成为争端；决策过程和张觉的地方基础并不完整，但双方围绕降将与缓冲地带的争执确为1125年金军南侵的直接借口之一。联盟曾承诺的燕云，在此已由目标变成一条难以供给、难以解释、也难以防守的边界。
